% !TEX program = xelatex
% =====================================================================
% LLM Theory Seminar Week 7: Chain-of-Thought & Reasoning Complexity
%
% Compile:
%   xelatex chain_of_thought_reasoning_complexity_notes.tex
%   xelatex chain_of_thought_reasoning_complexity_notes.tex
% or
%   latexmk -xelatex chain_of_thought_reasoning_complexity_notes.tex
%
% Korean typesetting: kotex + XeLaTeX.
%
% Fixed layout reference / inspiration:
%   - Ben Jones, "lecture-notes" (Overleaf, CC BY 4.0)
%   - CMU-Notes/notes-template (structured university course notes)
% This file intentionally keeps the approved seminar-note layout.
% =====================================================================

\documentclass[11pt,a4paper,oneside,openany]{memoir}

\usepackage{kotex}
\usepackage{amsmath,amssymb,amsthm,mathtools,bm}
\usepackage{booktabs,longtable,array,tabularx,makecell}
\newcolumntype{Y}{>{\raggedright\arraybackslash}X}
\usepackage{enumitem}
\usepackage[most]{tcolorbox}
\usepackage[dvipsnames]{xcolor}
\usepackage{xurl}
\usepackage[unicode]{hyperref}
\usepackage{cleveref}
\usepackage{needspace}

% ---------- page geometry ----------
\setlrmarginsandblock{27mm}{27mm}{*}
\setulmarginsandblock{24mm}{28mm}{*}
\checkandfixthelayout

\setlength{\parindent}{0pt}
\setlength{\parskip}{0.48em}
\setlength{\emergencystretch}{2em}
\setlength{\headheight}{15pt}
\linespread{1.055}\selectfont
\setlist[itemize]{topsep=0.28em,itemsep=0.12em,parsep=0pt,leftmargin=1.55em}
\setlist[enumerate]{topsep=0.28em,itemsep=0.16em,parsep=0pt,leftmargin=1.75em}

% ---------- restrained lecture-note palette ----------
\definecolor{NoteBlue}{HTML}{294E6B}
\definecolor{NoteBlueLight}{HTML}{F2F6F9}
\definecolor{NoteGray}{HTML}{5A6268}
\definecolor{NoteGrayLight}{HTML}{F6F6F4}
\definecolor{NoteWarm}{HTML}{7A6545}
\definecolor{NoteWarmLight}{HTML}{FAF7F0}

\hypersetup{
  colorlinks=true,
  linkcolor=NoteBlue,
  citecolor=NoteBlue,
  urlcolor=NoteBlue,
  pdftitle={Chain-of-Thought and Reasoning Complexity},
  pdfauthor={LLM Theory Seminar}
}

% ---------- running heads ----------
\makepagestyle{seminar}
\makeoddhead{seminar}{\footnotesize LLM Theory Seminar}{ }{\footnotesize Chain-of-Thought \& Reasoning Complexity}
\makeoddfoot{seminar}{}{\footnotesize\thepage}{}
\makeheadrule{seminar}{\textwidth}{0.35pt}
\pagestyle{seminar}
\aliaspagestyle{chapter}{seminar}

% ---------- chapter / section hierarchy ----------
\makechapterstyle{seminarnotes}{
  \setlength{\beforechapskip}{8pt}
  \setlength{\midchapskip}{8pt}
  \setlength{\afterchapskip}{22pt}
  \renewcommand{\chapterheadstart}{\vspace*{\beforechapskip}}
  \renewcommand{\chapnamefont}{\normalfont\small\bfseries}
  \renewcommand{\chapnumfont}{\normalfont\bfseries\fontsize{34}{38}\selectfont\color{NoteBlue}}
  \renewcommand{\chaptitlefont}{\normalfont\huge\bfseries\color{black}}
  \renewcommand{\printchaptername}{}
  \renewcommand{\chapternamenum}{}
  \renewcommand{\printchapternum}{\chapnumfont\thechapter}
  \renewcommand{\afterchapternum}{\par\nobreak\color{black}\vskip\midchapskip}
  \renewcommand{\printchaptertitle}[1]{\chaptitlefont ##1}
  \renewcommand{\afterchaptertitle}{\par\nobreak\vskip\afterchapskip}
}
\chapterstyle{seminarnotes}
\setsecheadstyle{\Large\bfseries}
\setsubsecheadstyle{\large\bfseries}
\setsubsubsecheadstyle{\normalsize\bfseries}
\setbeforesecskip{-1.3\baselineskip plus -0.2\baselineskip minus -0.1\baselineskip}
\setaftersecskip{0.55\baselineskip}
\setbeforesubsecskip{-0.9\baselineskip plus -0.15\baselineskip minus -0.1\baselineskip}
\setaftersubsecskip{0.35\baselineskip}

% ---------- notation ----------
\newcommand{\R}{\mathbb{R}}
\newcommand{\N}{\mathbb{N}}
\newcommand{\E}{\mathbb{E}}
\newcommand{\Pp}{\mathbb{P}}
\newcommand{\F}{\mathcal{F}}
\newcommand{\Hc}{\mathcal{H}}
\newcommand{\Lcal}{\mathcal{L}}
\newcommand{\D}{\mathcal{D}}
\newcommand{\CoT}{\mathsf{CoT}}
\newcommand{\TIME}{\mathsf{TIME}}
\newcommand{\SPACE}{\mathsf{SPACE}}
\newcommand{\VC}{\mathrm{VCdim}}
\newcommand{\Ldim}{\mathrm{Ldim}}
\newcommand{\Risk}{\mathcal{R}}
\newcommand{\Icot}{\mathcal{I}^{\mathrm{CoT}}}
\newcommand{\norm}[1]{\left\lVert #1\right\rVert}
\newcommand{\abs}[1]{\left\lvert #1\right\rvert}
\newcommand{\one}{\mathbf{1}}
\DeclareMathOperator*{\argmin}{arg\,min}
\DeclareMathOperator*{\argmax}{arg\,max}

% ---------- note environments ----------
\tcbset{
  seminarbox/.style={
    enhanced,breakable,sharp corners,boxrule=0pt,frame hidden,
    left=3.2mm,right=3mm,top=1.8mm,bottom=1.8mm,
    before={\par\Needspace{5\baselineskip}},before skip=0.8em,after skip=0.8em,
    fonttitle=\bfseries\small,coltitle=black,
    attach title to upper={\quad},
  }
}
\newtcolorbox{keybox}[1][]{
  seminarbox,colback=NoteBlueLight,
  borderline west={1.8pt}{0pt}{NoteBlue},
  title={핵심#1}
}
\newtcolorbox{paperbox}[1][]{
  seminarbox,colback=NoteGrayLight,
  borderline west={1.8pt}{0pt}{NoteGray},
  title={원 논문 기준#1}
}
\newtcolorbox{suppbox}[1][]{
  seminarbox,colback=NoteWarmLight,
  borderline west={1.8pt}{0pt}{NoteWarm},
  title={보충 유도 --- 논문 밖의 독자용 전개#1}
}
\newtcolorbox{warningbox}[1][]{
  seminarbox,colback=NoteGrayLight,
  borderline west={1.8pt}{0pt}{NoteGray},
  title={주의#1}
}
\newtcolorbox{presentbox}{
  seminarbox,colback=white,
  borderline west={2.2pt}{0pt}{black!70},
  fontupper=\footnotesize,
  top=1.2mm,bottom=1.2mm,before skip=0.5em,after skip=0.5em,
  title={발표에서 이렇게 말하기}
}

\theoremstyle{definition}
\newtheorem{definition}{정의}[chapter]
\newtheorem{proposition}{명제}[chapter]
\newtheorem{theorem}{정리}[chapter]
\newtheorem{lemma}{보조정리}[chapter]

\begin{document}

\begin{titlingpage}
\thispagestyle{empty}
\vspace*{1.2cm}
{\small\bfseries LLM Theory Seminar \hfill Week 7\par}
\vspace{1.4cm}
{\huge\bfseries Chain-of-Thought \& Reasoning Complexity\par}
\vspace{0.55cm}
{\large computational depth, in-context optimization, learning signals, reasoning-token lower bounds\par}
\vspace{0.9cm}
\rule{\textwidth}{0.7pt}
\vspace{1.1cm}

\begin{minipage}{0.86\textwidth}
\small
이번 주의 질문은 ``Chain-of-Thought(CoT)가 왜 잘 되는가?''를 하나의 메커니즘으로 설명하는 것이 아니다.
서로 다른 네 축을 분리한다. 첫째, autoregressive decoding이 추가하는 \textbf{순차 계산 깊이}. 둘째, Transformer의 한 forward pass 자체가 gradient descent나 backpropagation 같은 \textbf{학습 알고리즘을 내부에서 실행할 수 있는가}. 셋째, intermediate trace가 training에서 주는 \textbf{학습 신호와 sample-complexity 이득}. 넷째, inference에 필요한 \textbf{reasoning-token 예산과 하한}이다.
특히 이번 수정본에서는 두 번째 축 --- ICL과 optimization/backprop 사이의 구성적 동치성 --- 을 다른 세 축과 같은 비중으로 다룬다.
핵심은 ``길게 생각하면 강해진다''가 아니라, \textbf{forward pass 내부의 optimization과 autoregressive reasoning을 구분하고, 각각 무엇을 보장하는지}를 정확히 나누는 것이다.
\end{minipage}

\vfill
{\small
\textbf{Primary readings}\\[0.25em]
William Merrill and Ashish Sabharwal,
\textbf{The Expressive Power of Transformers with Chain of Thought}, ICLR 2024.\\
Nirmit Joshi et al.,
\textbf{A Theory of Learning with Autoregressive Chain of Thought}, COLT 2025.\\
Awni Altabaa, Omar Montasser, and John Lafferty,
\textbf{CoT Information: Improved Sample Complexity under Chain-of-Thought Supervision}, NeurIPS 2025.\\
Steve Hanneke, Idan Mehalel, and Shay Moran,
\textbf{Sample Complexity of Autoregressive Reasoning: Chain-of-Thought vs. End-to-End}, 2026.\\
Kiran Tomlinson et al.,
\textbf{Reasoning about Reasoning: BAPO Bounds on Chain-of-Thought Token Complexity in LLMs}, ICML 2026.\\
Yu Bai et al.,
\textbf{Transformers as Statisticians: Provable In-Context Learning with In-Context Algorithm Selection}, 2023.\\
Weimin Wu et al.,
\textbf{In-Context Deep Learning via Transformer Models}, ICML 2025.
\par}
\vspace{0.8cm}
{\small September 2026\par}
\end{titlingpage}

\tableofcontents*
\clearpage

% =====================================================================
\chapter{기호와 네 가지 이론 축}
% =====================================================================

\section{논문별 기호 대응표}

CoT 이론 논문은 같은 ``길이''를 서로 다른 대상으로 쓴다.
이번 노트에서는 input length를 $n$, generation length를 $T$, 실제 reasoning-token budget을 $R$로 구분한다.

\begin{center}
\small
\begin{tabularx}{\textwidth}{>{\bfseries\raggedright\arraybackslash}p{0.18\textwidth}Y Y Y}
\toprule
개념 & 이 노트 & Merrill--Sabharwal & Joshi / Hanneke / Altabaa \\
\midrule
입력 & $x\in\Sigma^*$, $|x|=n$ & input string $x$ & prompt $x$ \\
기본 next-token rule & $f:\Sigma^*\to\Sigma$ & fixed decoder block & base class $\F\subseteq\Sigma^{\Sigma^*}$ \\
중간 추론 & $z_{1:T}$ & decoding steps / scratchpad & full CoT trajectory \\
최종 출력 & $y$ & last generated answer token & $f^{\mathrm{e2e},T}(x)$ \\
CoT 전체 & $f^{\mathrm{CoT},T}(x)$ & generated sequence & observed label under CoT supervision \\
복잡도 자원 & $T$ or $R$ & decoding-step budget $t(n)$ & sample size $m$, $\VC(\F)$ \\
오차 & $\epsilon$ & language decision error보다 expressivity 중심 & PAC target error / end-to-end risk \\
\bottomrule
\end{tabularx}
\end{center}

\section{CoT를 함수 반복으로 쓰기}

고정된 next-token generator
\[
f:\Sigma^*\to\Sigma
\]
가 있다고 하자. 입력 $x$에서 autoregressive하게
\[
z_1=f(x),
\qquad
z_2=f(xz_1),
\qquad
\ldots,
\qquad
z_T=f(xz_{1:T-1})
\]
를 생성한다.

따라서 전체 chain과 최종 answer는
\[
f^{\mathrm{CoT},T}(x)
= z_{1:T},
\qquad
f^{\mathrm{e2e},T}(x)
= z_T
\]
로 구분할 수 있다.

여기서 중요한 점은 모든 시점에 \textbf{같은} $f$가 쓰인다는 것이다.
이 time-invariance가 이후 sample complexity 결과의 핵심 가정이다.

\section{네 가지 질문을 분리해야 한다}

\begin{keybox}
CoT 이론에서 서로 다른 네 질문을 섞지 않는다.
\begin{enumerate}
\item \textbf{Computational expressivity:} $T$개의 decoding step을 허용하면 어떤 함수를 계산할 수 있는가?
\item \textbf{In-context optimization:} 한 번의 Transformer forward pass 안에 GD나 backprop 같은 학습 알고리즘을 얼마나 충실하게 구현할 수 있는가?
\item \textbf{Statistical learnability:} 그런 autoregressive rule을 학습하려면 training example이 몇 개 필요한가?
\item \textbf{Inference token complexity:} 이미 rule을 알고 있을 때 문제를 풀기 위해 reasoning token이 최소 몇 개 필요한가?
\end{enumerate}
첫째는 계산복잡도, 둘째는 forward-pass 내부의 algorithmic simulation, 셋째는 sample complexity, 넷째는 inference-time lower bound다.
\end{keybox}

\begin{presentbox}
이번 주에는 CoT를 ``추론 문장을 길게 쓰는 기법''이 아니라 자원 관점으로 봅니다. Decoding step은 serial compute를 추가하고, forward-pass depth는 context 안에서 optimization loop를 펼칠 수 있게 하며, intermediate supervision은 학습 신호를 추가하고, reasoning token budget은 inference 비용을 제한합니다. 이 네 효과는 수학적으로 서로 다른 문제입니다.
\end{presentbox}

% =====================================================================
\chapter{CoT는 순차 계산 자원이다}
% =====================================================================

\section{Week 5와의 연결: 한 번의 forward pass와 반복 decoding}

Week 5에서 log-precision, fixed-depth Transformer의 한 번의 계산을 강하게 병렬화 가능한 회로로 보는 관점을 다뤘다.
Week 7에서는 그 결과를 반복하지 않고, \textbf{decoding step을 추가하면 무엇이 달라지는가}에만 초점을 둔다.

Decoder가 한 step마다
\[
y_{t+1}=F_\theta(x,y_{1:t})
\]
를 수행하면, 모델 depth가 고정되어 있어도 $t=1,2,\ldots,T$에 걸쳐 recurrent computation이 생긴다.
따라서 총 sequential depth는 대략
\[
\text{effective sequential depth}
\propto T
\]
로 증가한다.

\section{DFA simulation이 보여주는 가장 단순한 이유}

Deterministic finite automaton을
\[
A=(Q,\Sigma,\delta,q_0,F)
\]
라고 하자. 문자열 $x_1\cdots x_n$을 읽을 때 상태는
\[
q_i=\delta(q_{i-1},x_i)
\]
로 갱신된다.

이 식은 본질적으로 recurrence다. 만약 reasoning token $z_i$가 $q_i$를 저장한다면
\[
z_i=q_i
\]
로 둘 수 있고, $i$번째 decoding step에서
\[
(q_{i-1},x_i)
\mapsto q_i
\]
를 한 번 계산하면 된다.

Merrill--Sabharwal의 Theorem 1은 적절한 projected pre-norm, saturated attention 가정 아래 모든 regular language를 $|x|+1$ decoding steps로 인식하는 decoder-only Transformer를 구성한다 \cite{merrill2024}.

\begin{paperbox}
이 결과는 ``일반적인 상용 Transformer가 regular language를 항상 학습한다''는 뜻이 아니다.
존재성 construction이며, strict causal saturated attention과 projected pre-norm이라는 명시적 모델 가정이 있다.
핵심 이론 메시지는 \textbf{linear number of decoding steps가 recurrence를 가능하게 한다}는 것이다.
\end{paperbox}

\section{Turing-machine simulation과 계산 하한/상한}

Merrill--Sabharwal의 Theorem 2는 polynomially bounded time $t(n)$ 동안 동작하는 Turing machine $M$에 대해, 한 decoding step으로 한 TM step을 simulate하는 Transformer construction을 준다.
따라서 논문 표현으로
\[
\TIME(t(n))\subseteq\CoT(t(n))
\]
이다.

반대 방향에서는 계산을 simulate하는 비용을 세어
\[
\CoT(t(n))
\subseteq
\widetilde{\TIME}\!\left(n^2+t(n)^2\right)
\]
를 얻고, 공간 측면에서는
\[
\CoT(t(n))
\subseteq
\SPACE\!\left(t(n)+\log n\right)
\]
를 얻는다 \cite{merrill2024}.

\begin{suppbox}
왜 space upper bound에 $t(n)$이 직접 들어가는지 직관적으로 확인해 보자.
생성한 intermediate token을 모두 buffer에 저장한다고 하면 필요한 저장량은
\[
O(t(n)).
\]
각 다음 token 계산을 위한 fixed-depth log-precision Transformer simulation에 logarithmic workspace가 필요하다고 두면
\[
O(\log(n+t(n))).
\]
따라서 총 workspace는
\[
O\!\left(t(n)+\log(n+t(n))\right)
=O\!\left(t(n)+\log n\right)
\]
이다. 마지막 등식은 논문의 polynomially bounded $t(n)$ 가정 아래 이해하면 된다.
\end{suppbox}

\section{왜 ``CoT = P''라고 단순화하면 안 되는가}

Polynomial-length CoT가 polynomial-time computation까지 도달한다는 요약은 유용하지만, 반드시 다음 조건을 붙여야 한다.

\begin{itemize}
\item 모델은 논문에서 정의한 특정 decoder architecture다.
\item precision, normalization, attention 방식에 대한 가정이 있다.
\item 결과는 \textbf{representability}다. gradient-based training이 그 construction을 찾는다는 정리가 아니다.
\item CoT length $T$는 계산 자원이며, 실제 latency와 KV-cache cost는 별도다.
\end{itemize}

\begin{presentbox}
CoT가 주는 첫 번째 효과는 serial depth입니다. 한 번의 forward pass가 병렬 회로라면, autoregressive decoding은 그 회로를 반복 호출하면서 이전 출력 토큰을 다음 계산의 상태로 씁니다. 그래서 reasoning-token 수가 단순한 문장 길이가 아니라 계산 시간에 가까운 자원이 됩니다.
\end{presentbox}

% =====================================================================
\chapter{ICL 안의 optimization loop: attention에서 backprop까지}
% =====================================================================

\section{출발점: linear attention 한 층은 GD 한 step을 구현할 수 있다}

먼저 ``ICL이 학습한다''는 말을 가장 좁고 검증 가능한 형태로 쓴다. Reference predictor를
\[
f_W(x)=Wx,
\qquad
L(W)=\frac{1}{2N}\sum_{i=1}^N\norm{Wx_i-y_i}^2
\]
라고 하자. 그러면
\[
\nabla_WL(W)
=\frac1N\sum_{i=1}^N(Wx_i-y_i)x_i^\top
\]
이고, learning rate $\eta$의 한 gradient-descent step은
\[
\Delta W
=-\eta\nabla_WL(W)
=-\frac{\eta}{N}\sum_{i=1}^N(Wx_i-y_i)x_i^\top.
\]
Query $x_j$에서 prediction이 바뀌는 양은
\[
\Delta W x_j
=-\frac{\eta}{N}\sum_{i=1}^N
(Wx_i-y_i)(x_i^\top x_j).
\]
즉 각 context example의 residual과 query와의 similarity를 곱해 합한다.

von Oswald et al.의 핵심 observation은 linear self-attention에도 정확히 같은 bilinear sum이 있다는 것이다 \cite{vonoswald2023}. 한 head를
\[
e_j^+=e_j+PVK^\top q_j
=e_j+P\sum_{i=1}^Nv_i(k_i^\top q_j)
\]
로 쓰고 token을
\[
e_i=\begin{bmatrix}x_i\\y_i\end{bmatrix}
\]
로 둔다. 다음 block projections을 택하자.
\[
W_K=W_Q=
\begin{bmatrix}I&0\\0&0\end{bmatrix},
\qquad
W_V=
\begin{bmatrix}0&0\\W&-I\end{bmatrix}.
\]
그러면
\[
k_i^\top q_j=x_i^\top x_j,
\qquad
v_i=\begin{bmatrix}0\\Wx_i-y_i\end{bmatrix}.
\]
Output projection의 $y$ block을 $\eta I/N$로 잡으면 attention이 만드는 $y$-coordinate 변화는
\[
\frac{\eta}{N}\sum_{i=1}^N(Wx_i-y_i)(x_i^\top x_j)
=-\Delta W x_j.
\]
따라서
\[
\boxed{
 e_j^+
 =
 \begin{bmatrix}
 x_j\\y_j-\Delta W x_j
 \end{bmatrix}}
\]
가 된다. 특히 query $x_q$의 $y$-slot을 $-Wx_q$로 초기화하고 마지막에 부호를 뒤집어 readout하면
\[
-(-Wx_q-\Delta W x_q)=(W+\Delta W)x_q
\]
이므로 한 GD step 뒤의 prediction을 정확히 얻는다. 이 exact construction에서는 context token만 key/value에 넣는 설정(또는 $W\approx0$인 단순화)도 함께 사용된다.

\begin{paperbox}
von Oswald et al. Proposition 1의 정확한 범위는 \textbf{one-head linear self-attention + 특정 regression tokenization}이다. Softmax를 제거한 attention에 대해 한 층의 data transformation과 MSE의 GD 한 step이 정확히 같아지도록 하는 parameter construction이 존재한다. 논문은 추가로 단순 regression setting에서 실제 meta-training이 이 construction에 가까운 해를 찾는 경우를 실험한다.
\end{paperbox}

Aky\"urek et al.은 같은 문제를 더 일반적인 Transformer decoder primitive로 다룬다 \cite{akyurek2023}. Theorem 1은 한 in-context example에 대한 gradient update와 prediction을 constant number of layers, $O(d)$ hidden space로 구현할 수 있음을 construction으로 보인다. 여러 update는 이 block을 쌓아 구현한다. 따라서 초기 연구의 핵심 명제는
\[
\text{context} \longmapsto \text{task-specific optimizer state}
\]
를 forward activations 안에 기록할 수 있다는 것이다.

\section{Bai et al.: GD를 base ICL algorithm으로 만들고 algorithm까지 선택한다}

Bai et al.은 ``한 번의 GD step을 흉내 낸다''에서 두 단계 더 나아간다 \cite{bai2023}.
첫째, smooth convex empirical risk
\[
\widehat L_N(w)
=\frac1N\sum_{i=1}^N\ell(w^\top x_i,y_i)
\]
에 대해
\[
w^{t+1}=w^t-\eta\nabla\widehat L_N(w^t)
\]
를 context 안에서 반복 실행하는 construction을 준다. Theorem 13의 핵심 형태는 \textbf{$L$번의 GD iterate를 $(L+1)$-layer attention-only Transformer가 근사}한다는 것이다. 한 attention layer가 한 in-context GD iterate를 구현하고, convexity/smoothness를 이용해 근사오차가 $L$에 대해 선형으로 누적되도록 제어한다.

여기서 중요한 세부사항은 이론 construction이 표준 softmax가 아니라 \textbf{normalized ReLU attention}을 사용한다는 점이다. 즉
\[
\text{``Transformer가 GD를 한다''}
\]
보다
\[
\text{``특정 attention primitive가 GD iterate를 구현할 수 있다''}
\]
가 정확한 문장이다.

둘째, Bai et al.은 base algorithm 여러 개를 같은 Transformer가 context에 따라 선택하게 한다. 개념적으로 post-ICL validation은
\[
\widehat k
=\argmin_{k\in[K]}
\widehat R_{\mathrm{val}}
\bigl(A_k(D_{\mathrm{train}})\bigr)
\]
처럼 후보 ICL algorithm $A_1,\ldots,A_K$를 실행한 뒤 validation signal로 고른다. 반대로 pre-ICL testing은 먼저 context statistic을 계산해
\[
\widehat k=g(s(D))
\]
를 정하고 그 algorithm을 실행한다. 논문은 두 routing mechanism 모두에 explicit transformer construction을 준다.

\begin{keybox}
이 결과가 주는 관점 변화는 ``Transformer가 하나의 optimizer를 흉내 낸다''에서 ``Transformer가 context 안에서 \textbf{어떤 learning algorithm을 실행할지까지 선택할 수 있다}''로의 확장이다. ICL을 단일 estimator가 아니라 algorithm-selection problem으로 볼 수 있다.
\end{keybox}

\section{Wu et al.: forward pass 안에서 깊은 network의 backprop을 시뮬레이션한다}

Linear regression의 한 step과 깊은 network의 backprop 사이에는 큰 간격이 있다. Wu et al.은 이 간격을 explicit construction으로 메운다 \cite{wu2025}. $N$-layer network의 parameter를 $w$라 하고 empirical risk를 $L_n(w)$라 하자. Target inner optimizer는
\[
w^{(\ell+1)}
=\operatorname{Proj}_{\mathcal W}
\left(w^{(\ell)}-\eta\nabla L_n(w^{(\ell)})\right).
\]

한 GD step을 계산하려면 먼저 network의 forward propagation으로 모든 activation을 만들고, 그 뒤 chain rule을 거꾸로 적용해 각 layer의 gradient를 계산해야 한다. 논문의 construction은 이 계산을 Transformer depth에 펼친다.

\begin{enumerate}
\item 약 $N$개 layer가 in-context example들의 forward activations를 계산한다.
\item 다음 약 $N$개 layer가 backward error/chain-rule 항을 역순으로 전달한다.
\item 추가 layer들이 sample-wise gradient를 집계하고 parameter update와 projection을 수행한다.
\end{enumerate}

Theorem 1의 정량적 statement는 다음 형태다.
\[
\boxed{
\text{$L$ GD steps on an $N$-layer network}
\quad\Longleftarrow\quad
(2N+4)L\text{-layer Transformer}}
\]
그리고 각 simulated step은
\[
\widehat w^{(\ell)}
=
\operatorname{Proj}_{\mathcal W}
\left(
\widehat w^{(\ell-1)}
-\eta\nabla L_n(\widehat w^{(\ell-1)})
+\varepsilon^{(\ell-1)}
\right),
\qquad
\norm{\varepsilon^{(\ell-1)}}_2\le \eta\epsilon
\]
형태의 inexact GD update를 구현한다. Primary construction은 ReLU-attention 계열을 쓰며, 논문은 softmax Transformer로의 extension도 별도로 분석한다.

\begin{warningbox}
arXiv:2411.16549의 현재 제목은 In-Context Deep Learning via Transformer Models이다. ``How Do Transformers Learn In-Context Beyond Simple Functions?''는 arXiv:2310.10616의 별도 논문이다. $(2N+4)L$ backprop simulation theorem은 arXiv:2411.16549의 Wu et al. 결과이므로 이 노트에서는 해당 원문을 기준으로 쓴다. 논문 본문은 ``ReLU neural network''라는 표현을 쓰지만, Theorem 1의 formal assumptions에는 target activation $r,r'$의 smoothness가 들어간다. 따라서 nonsmooth ReLU 자체에 대한 무조건적인 exact theorem으로 읽으면 안 된다.
\end{warningbox}

\section{encoder냐 decoder냐보다 중요한 두 축}

이 계열의 논문을 architecture 이름만으로 분류하면 본질을 놓치기 쉽다.

\begin{center}
\small
\begin{tabularx}{\textwidth}{>{\bfseries\raggedright\arraybackslash}p{0.18\textwidth}YYY}
\toprule
논문 & attention / architecture & 내부에서 구현하는 것 & 핵심 한계 \\
\midrule
von Oswald et al. & linear self-attention & linear regression GD 1 step & exact construction이 softmax가 아님 \\
Aky\"urek et al. & decoder construction & GD / ridge-related updates & linear-model setting 중심 \\
Bai et al. & normalized ReLU attention; encoder 기본, decoder extension & multi-step GD + algorithm selection & 주로 explicit existence construction \\
Wu et al. & ReLU construction + softmax extension & deep network의 forward/backprop + GD & depth가 $(2N+4)L$로 증가 \\
\bottomrule
\end{tabularx}
\end{center}

따라서 더 중요한 분류축은 다음 둘이다.

\begin{enumerate}
\item \textbf{linear / ReLU-style attention에서 softmax attention으로 얼마나 가까이 가는가};
\item \textbf{얕은 estimator update를 구현하는가, 아니면 깊은 network의 forward/backward sweep 전체를 구현하는가}.
\end{enumerate}

Encoder/decoder 차이는 causal mask와 information routing에 영향을 주지만, 이 theorem family의 중심 질문은 ``forward computation 안에 어떤 optimizer circuit을 심을 수 있는가''이다.

\section{Chapter 2와 faithfulness에 어떻게 연결되는가}

Chapter 2의 CoT 결과는 \textbf{여러 번의 autoregressive forward pass}를 이어 붙여 serial computation을 만든다.
반면 지금 본 결과는 \textbf{한 번의 forward pass 내부의 layer depth}에 optimization loop를 펼친다. 두 자원을 구분해
\[
D_{\mathrm{inner}}
=\text{optimizer/backprop을 구현하는 network depth},
\]
\[
T_{\mathrm{decode}}
=\text{autoregressive reasoning steps}.
\]
라고 생각하는 편이 안전하다. 둘은 동시에 존재할 수도 있다.

또 ``in-context으로 학습한다''는 말에도 세 수준이 있다.

\begin{enumerate}
\item \textbf{Representability}: 그런 optimizer와 같은 함수를 계산하는 Transformer parameter가 존재한다.
\item \textbf{Learnability of the circuit}: outer-loop SGD가 그 parameter/circuit을 실제로 찾는다.
\item \textbf{Mechanistic faithfulness}: 실제 자연 task에서 내부 computation이 그 optimizer라는 설명과 인과적으로 일치한다.
\end{enumerate}

위 construction theorem들이 가장 강하게 보장하는 것은 1번이다. von Oswald와 Aky\"urek은 2번에 대한 실험적 evidence도 제공하지만, 1번에서 2번이나 3번이 자동으로 따라오지는 않는다. 이 점이 뒤의 faithfulness 절과 직접 연결된다.

\begin{presentbox}
ICL--backprop 동치성 축에서 중요한 건 encoder냐 decoder냐가 아닙니다. 한쪽 끝에는 linear attention 한 층이 regression GD 한 step을 정확히 구현하는 결과가 있고, 다른 끝에는 $(2N+4)L$층 Transformer가 $N$층 network의 backprop을 $L$번 시뮬레이션하는 결과가 있습니다. 즉 forward pass 자체가 작은 training loop를 품을 수 있다는 것이 핵심입니다. 다만 대부분은 ``그런 circuit이 존재한다''는 정리이지, SGD가 실제 LLM에서 항상 그 circuit을 찾는다는 정리는 아닙니다.
\end{presentbox}

% =====================================================================
\chapter{Autoregressive CoT를 학습한다는 것}
% =====================================================================

\section{Joshi et al.의 formal learning model}

Joshi et al.은 base next-token class
\[
\F\subseteq \Sigma^{\Sigma^*}
\]
를 두고, 한 $f\in\F$를 $T$번 반복하여 얻는 end-to-end class를 연구한다 \cite{joshi2025}.

Training supervision은 두 가지다.

\begin{enumerate}
\item \textbf{E2E supervision}: $(x,z_T)$만 관측한다.
\item \textbf{CoT supervision}: $(x,z_1,\ldots,z_T)$ 전체를 관측한다.
\end{enumerate}

같은 latent generator $f_\star$가 두 setting 모두의 ground truth다.
차이는 learner가 intermediate transition을 직접 보느냐이다.

\section{왜 E2E에서는 \texorpdfstring{$T$}{T}가 어려움을 키울 수 있는가}

$T$-step map을
\[
\F^{\mathrm{e2e},T}
=\left\{f^{\mathrm{e2e},T}:f\in\F\right\}
\]
라고 하자.
Base class가 단순하더라도 $T$번 composition한 최종 mapping은 더 복잡해질 수 있다.

Joshi et al.은 binary alphabet에서 대략
\[
\VC\!\left(\F^{\mathrm{e2e},T}\right)
\le 6T\,\VC(\F)
\]
형태의 bound를 보인다.
표준 realizable PAC bound를 쓰면
\[
m_{\mathrm{e2e}}
=
O\!\left(
\frac{T\,\VC(\F)\log(1/\epsilon)+\log(1/\delta)}{\epsilon}
\right)
\]
가 나온다.

\begin{paperbox}
더 중요한 것은 이 linear-$T$ dependence가 단순 proof artifact만은 아니라는 점이다.
Joshi et al. Theorem 3.3은 주어진 VC dimension $D$와 generation length $T$에 대해, E2E learning이 $\Omega(DT)$ 규모의 sample dependence를 실제로 요구할 수 있는 base class를 구성한다.
\end{paperbox}

\section{CoT supervision을 보면 무엇이 달라지는가}

한 trajectory
\[
z=(x,z_1,\ldots,z_T)
\]
를 관측했다고 하자.
그러면 사실상 같은 base rule $f_\star$에 대한 transition constraint를 $T$개 얻는다.

\[
f_\star(x)=z_1,
\]
\[
f_\star(xz_1)=z_2,
\]
\[
\cdots
\]
\[
f_\star(xz_{1:T-1})=z_T.
\]

즉 하나의 trajectory는 final label 하나보다 훨씬 구조화된 정보를 준다.
Joshi et al.의 CoT consistency rule은 training trajectory 전체와 일치하는 $\hat f\in\F$를 찾는다.

그 결과 binary case에서
\[
m_{\mathrm{CoT}}
=
O\!\left(
\frac{
\VC(\F)\log T\,\log(1/\epsilon)
+
\log(1/\delta)
}{\epsilon}
\right)
\]
형태의 bound를 얻는다.
E2E의 최악 linear-$T$ dependence가 log-$T$로 줄어든다.

\section{Littlestone dimension이 등장하는 이유}

E2E learning에서도 stronger combinatorial assumption을 두면 $T$ dependence를 줄일 수 있다.
Joshi et al.은
\[
\VC\!\left(\F^{\mathrm{e2e},T}\right)
\lesssim
\Ldim(\F)\log T
\]
형태의 결과를 보인다.
따라서
\[
m_{\mathrm{e2e}}
=
\widetilde O\!\left(
\frac{\Ldim(\F)\log T}{\epsilon}
\right)
\]
가 가능하다.

왜 online-learning dimension이 나오는가? Autoregressive rollout에서는 learner가 이전 generated prefix에 따라 다음 query point를 스스로 만든다.
즉 fixed i.i.d. input classification보다 tree-structured sequential interaction에 가깝다.

\begin{warningbox}
``CoT를 보면 sample complexity가 $T$에 독립이다''는 문장은 Joshi et al. 2025 결과만 인용하면 정확하지 않다.
그 논문의 일반 VC-based CoT bound에는 $\log T$가 남는다.
완전한 $T$-independence는 이후 Hanneke--Mehalel--Moran 2026이 더 정교한 learning rule로 개선한다.
\end{warningbox}

\begin{presentbox}
CoT supervision의 통계적 이점은 중간 문장이 친절해서가 아닙니다. 같은 time-invariant next-token rule에 대해 한 trajectory가 여러 transition constraint를 동시에 줍니다. 그래서 final answer만 보는 것보다 hypothesis를 훨씬 빨리 제거할 수 있습니다.
\end{presentbox}

% =====================================================================
\chapter{2026 결과: CoT sample complexity에서 \texorpdfstring{$T$}{T}를 없애기}
% =====================================================================

\section{핵심 정리와 직관}

Hanneke--Mehalel--Moran은 Joshi et al.의 일반 VC-based CoT bound에 남아 있던 $\log T$가 필연적인지 묻는다 \cite{hanneke2026}. Finite-VC base class에서는 적절한 learner를 쓰면 \textbf{CoT supervision의 sample complexity에서 generation length $T$를 제거할 수 있다}. 대표 형태는
\[
m^{\mathrm{CoT},T}_{\F}(\epsilon,\delta)
\le c(\F)
\left(
\frac1\epsilon\log\frac1\epsilon
+\log\frac1\delta
\right),
\]
여기서 $c(\F)$는 base class의 combinatorial complexity에는 의존하지만 $T$에는 의존하지 않는다.

직관은 한 trajectory
\[
(x;z_1,\ldots,z_T)
\]
가 동일한 time-invariant rule $f_\star$에 대한 transition constraints
\[
(x,z_1),\ (xz_1,z_2),\ldots,(xz_{1:T-1},z_T)
\]
를 제공한다는 것이다. 이 transition들은 iid sample은 아니지만, compression 관점에서는 hypothesis를 결정하는 핵심 support set의 크기가 trajectory length와 함께 커질 필요가 없다.

\section{E2E와 CoT를 섞지 말아야 한다}

같은 결과가 end-to-end supervision에도 자동 적용되는 것은 아니다. E2E에서는 intermediate transition을 보지 못하므로 class 구조에 따라 $T$ dependence가 constant에서 linear까지 다양하게 나타날 수 있다.

\begin{keybox}
Inference에서 CoT를 생성하는 것은 \textbf{computational resource}를 추가하는 일이고, training에서 CoT trace를 label로 주는 것은 \textbf{statistical supervision}을 추가하는 일이다. Hanneke et al.의 $T$-independence는 두 번째 의미의 결과다. Token 수나 training FLOPs가 공짜라는 뜻은 아니다.
\end{keybox}

\begin{presentbox}
2025 결과에서는 CoT supervision bound에 $\log T$가 남았지만, 2026 후속 결과는 finite-VC base class에서 적절한 learner를 쓰면 sample 수의 $T$ dependence 자체를 없앨 수 있음을 보입니다. 긴 trace가 더 많은 token 계산을 요구하는 것과, 더 많은 독립 training example을 요구하는 것은 다른 문제입니다.
\end{presentbox}

% =====================================================================
\chapter{CoT Information: 중간 과정이 얼마나 informative한가}
% =====================================================================

\section{Hypothesis discrimination으로 보는 CoT}

Altabaa--Montasser--Lafferty는 CoT supervision의 이득을 ``bad hypothesis를 얼마나 빨리 구분하는가''로 정량화한다 \cite{altabaa2025}. 두 hypothesis가 한 sample에서 CoT와 final answer 모두 같아 보일 확률을 $p_{\mathrm{agree}}$라 하면
\[
\Icot_{\D}(h_1,h_2)
=-\log p_{\mathrm{agree}}.
\]
Target $h_\star$에 대해 end-to-end error가 적어도 $\epsilon$인 후보만 모아
\[
\Icot_{\D,h_\star}(\epsilon;\Hc)
=
\inf_{\Risk^{\mathrm{e2e}}(h)\ge\epsilon}
\Icot_\D(h,h_\star)
\]
를 정의한다. 값이 클수록 틀린 hypothesis가 intermediate trace에서 빨리 드러난다.

\section{Sample complexity에서 어떻게 나타나는가}

Finite class에서 CoT-consistent hypothesis를 택하면 bad hypothesis 하나가 $m$개 sample 동안 target과 구분되지 않을 확률은 대략
\[
\exp\{-m\Icot_\D(h,h_\star)\}
\]
이다. Union bound로
\[
\boxed{
 m\gtrsim
 \frac{\log|\Hc|+\log(1/\delta)}
 {\Icot_{\D,h_\star}(\epsilon;\Hc)}
}
\]
를 얻는다. 또한 agreement 사건은 ``final answer도 같다''는 사건의 부분집합이므로
\[
\Icot_{\D,h_\star}(\epsilon;\Hc)\ge -\log(1-\epsilon)\ge\epsilon.
\]
따라서 realizable setting에서 informative한 CoT는 E2E signal보다 더 빠른 hypothesis elimination을 가능하게 할 수 있다.

\begin{warningbox}
$\Icot$은 설명의 자연스러움, factuality, causal faithfulness를 재는 quantity가 아니다. 오직 statistical distinguishability를 측정한다. ``중간 trace가 informative하다''와 ``중간 trace가 내부 계산을 정직하게 설명한다''는 별개다.
\end{warningbox}

\begin{presentbox}
CoT information은 reasoning 문장이 좋아 보이는지를 재는 점수가 아니라, 틀린 hypothesis를 training trace가 얼마나 빨리 배제하는지를 재는 통계적 양입니다. 그래서 sample complexity의 분모에 직접 들어갑니다.
\end{presentbox}

% =====================================================================
\chapter{Reasoning token 자체의 하한: BAPO}
% =====================================================================

\section{왜 token lower bound가 필요한가}

앞의 sample-complexity 결과는 training example 수를 다뤘다. Tomlinson et al.은 BAPO, 즉 bounded attention prefix oracle이라는 모형을 사용해 이미 학습된 reasoning system의 \textbf{inference-time token complexity}를 묻는다 \cite{tomlinson2026}. 한 step에서 prefix로 전달할 수 있는 summary와 직접 inspect할 수 있는 token 수를 제한하면, reasoning token 하나가 제한된 bandwidth의 communication round가 된다.

\section{핵심 lower bound}

Constant-bandwidth BAPO-CoT에서 논문은 대표적으로
\[
\mathrm{Majority}:\ \Omega(n),
\qquad
\mathrm{Match3}:\ \Omega(n),
\]
그리고 graph reachability에서는 edge 수 $m$에 대해
\[
\mathrm{Reachability}:\ \Omega(m)
\]
형태의 lower bound를 보인다.

Proof의 공통 골격은 fooling pair다. 정답이 다른 두 input $x,x'$를 만들되, $R$이 너무 작으면 제한된 bandwidth 때문에 각 reasoning step에서 관찰 가능한 정보를 같게 유지한다. 그러면 induction으로
\[
z_t(x)=z_t(x')\qquad (t\le R)
\]
가 되어 final answer도 같아지는 contradiction이 생긴다. 따라서 충분한 communication rounds, 즉 reasoning tokens가 필요하다.

\begin{warningbox}
BAPO theorem은 실제 GPT류 LLM 전체에 대한 lower bound가 아니다. 특정 bounded-information-flow abstraction에 대한 결과다. 또한 reachability의 정밀한 statement는 edge count $m$을 쓰므로 ``항상 $\Omega(n)$''으로 뭉뚱그리지 않는다.
\end{warningbox}

\section{학습 신호 축과의 대비}

Hanneke 결과에서 CoT sample complexity가 $T$에 독립일 수 있다는 사실과 BAPO의 $\Omega(n)$ token lower bound는 모순이 아니다.
\[
\underbrace{m}_{\text{training examples}}
\quad\text{와}\quad
\underbrace{R}_{\text{inference reasoning tokens}}
\]
는 다른 자원이다. 하나는 ``규칙을 몇 trajectory로 배울 수 있는가'', 다른 하나는 ``이미 배운 규칙으로 한 instance를 풀 때 몇 communication round가 필요한가''를 묻는다.

\begin{presentbox}
Reasoning token은 전부 불필요한 verbosity가 아닙니다. BAPO처럼 한 step의 information flow를 제한한 모델에서는 일부 문제에 선형 token 하한이 생깁니다. 따라서 training sample 수를 줄이는 것과 inference reasoning 길이를 줄이는 것은 전혀 다른 이론 문제입니다.
\end{presentbox}

% =====================================================================
\chapter{길게 생각하는 것의 비용과 오류 누적}
% =====================================================================

\section{필요한 reasoning length와 불필요한 verbosity는 다르다}

앞 장의 lower bound는 어떤 task에는 긴 reasoning이 필요하다는 뜻이다.
그러나 필요 이상으로 token을 늘리는 것이 항상 이득이라는 뜻은 아니다.

총 inference cost를 거칠게
\[
C_{\mathrm{total}}(R)
=
C_{\mathrm{prompt}}
+
\sum_{t=1}^{R}C_{\mathrm{decode}}(t)
\]
로 두자.
KV cache를 사용해도 각 step의 attention/read cost는 context length와 무관하지 않다.
따라서 $R$ 증가는 latency와 memory traffic을 늘린다.

\section{Local error가 독립이라는 toy model}

\begin{suppbox}
다음은 특정 논문의 정리가 아니라 오류 누적을 설명하기 위한 독자용 toy model이다.
각 reasoning step이 올바를 확률을 $1-p$라 하고, 모든 step이 맞아야 final answer가 맞으며 step error가 독립이라고 가정하자.
그러면 chain 전체 성공 확률은
\[
P_{\mathrm{succ}}(R)
=(1-p)^R.
\]
작은 $p$에 대해
\[
\log P_{\mathrm{succ}}(R)
=R\log(1-p)
\approx -pR,
\]
따라서
\[
P_{\mathrm{succ}}(R)
\approx e^{-pR}.
\]
즉 local reliability가 고정되어 있으면 chain이 길수록 누적 실패 확률은 커질 수 있다.
실제 reasoning model은 error correction, backtracking, self-verification 때문에 이 독립 모형과 다르다.
\end{suppbox}

\section{그래서 optimal reasoning budget 문제가 생긴다}

Task 난이도에 따라 필요한 최소 계산량을 $R_{\min}(x)$라고 하자.
그보다 짧으면 computational lower bound 때문에 실패할 수 있고,
너무 길면 비용과 오류 누적이 커질 수 있다.

개념적으로는
\[
R^*(x)
=
\argmin_R
\left[
\Risk_{\mathrm{task}}(R\mid x)
+\lambda C_{\mathrm{total}}(R)
\right]
\]
같은 문제로 볼 수 있다.

이 식은 이번 주 primary paper의 theorem이 아니라 inference-time compute allocation을 연결하기 위한 모델링 식이다.
그럼에도 중요한 메시지는 분명하다.

\begin{keybox}
``더 긴 CoT''와 ``더 많은 유효 계산''은 동일하지 않다.
필요한 serial depth를 확보하는 token과, 이미 충분한 계산 이후의 verbosity를 구분해야 한다.
이 distinction이 adaptive reasoning / early stopping / verifier 연구로 연결된다.
\end{keybox}

\section{Faithfulness 문제는 별도 축이다}

Visible CoT가 실제 내부 computation을 정확히 반영하는가도 별개의 질문이다. 이번 주의 complexity theorem들은 intermediate token을 computational state 또는 supervision signal로 취급하고, 앞의 ICL--optimization construction은 특정 forward function이 GD/backprop과 같아질 수 있음을 보였다. 둘 모두에서 \textbf{functional equivalence와 mechanistic explanation을 구분}해야 한다.

함수적 동치만으로 내부 메커니즘까지 결론내릴 수는 없다.
\[
\begin{aligned}
&\text{GD/backprop과 같은 output을 계산한다}\not\Rightarrow
\text{그 알고리즘이 실제 learned circuit이다},\\
&\text{CoT가 계산에 유용하다}\not\Rightarrow
\text{visible CoT가 내부 인과 과정을 설명한다}.
\end{aligned}
\]

따라서 ``ICL이 진짜 학습인가?''라는 질문에는 먼저 의미를 고정해야 한다. Context에 따라 output이 optimizer update와 함수적으로 동치라는 의미라면 construction theorem이 답할 수 있다. 반면 outer-loop SGD가 그 circuit을 찾았는지, visible rationale가 그 hidden computation을 설명하는지는 별도의 mechanistic/causal question이다.

\begin{presentbox}
Reasoning length에는 양면성이 있습니다. 너무 짧으면 필요한 serial computation을 못 하고, 너무 길면 latency와 오류 누적이 늘 수 있습니다. 그래서 중요한 질문은 ``CoT가 긴가 짧은가''가 아니라, task마다 필요한 최소 계산량과 실제 사용한 계산량의 gap입니다.
\end{presentbox}

% =====================================================================
\chapter{세 이론을 하나의 그림으로 연결하기}
% =====================================================================

\section{Inference와 training을 한 표에 놓기}

\begin{center}
\small
\begin{tabularx}{\textwidth}{>{\bfseries\raggedright\arraybackslash}p{0.21\textwidth}YYY}
\toprule
관점 & 질문 & 핵심 자원 & 대표 결과 \\
\midrule
Computational expressivity & CoT를 허용하면 무엇을 계산 가능한가? & decoding steps $T$ & $\TIME(t)\subseteq\CoT(t)$ under assumptions \\
ICL $\leftrightarrow$ optimization & 한 forward pass가 learning algorithm을 실행할 수 있는가? & layer depth, attention primitive & 1-layer GD부터 $(2N+4)L$ backprop simulation까지 \\
Autoregressive learnability & $T$-step rule을 몇 sample로 학습하는가? & examples $m$, $\VC(\F)$ & CoT supervision은 $T$ dependence를 줄이거나 제거 가능 \\
CoT information & trace가 bad hypothesis를 얼마나 잘 구별하는가? & $\Icot(\epsilon)$ & sample complexity $\sim d/\Icot(\epsilon)$ \\
Inference token lower bound & 최소 몇 reasoning token이 필요한가? & token budget $R$ & BAPO-hard tasks에서 linear lower bounds \\
\bottomrule
\end{tabularx}
\end{center}

\section{가장 중요한 다섯 개의 식}

Autoregressive recurrence:
\[
z_{t+1}=f(xz_{1:t}).
\]
Serial computation:
\[
\TIME(t(n))\subseteq\CoT(t(n)).
\]
In-context gradient step:
\[
\Delta W
=-\frac{\eta}{N}\sum_i(Wx_i-y_i)x_i^\top.
\]
Linear attention은 prediction에 같은 update를 구현할 수 있다.

Deep backprop simulation:
\[
\begin{gathered}
L\text{ GD steps on an }N\text{-layer network}\\
\Longleftarrow (2N+4)L\text{ Transformer layers}.
\end{gathered}
\]
Learning/inference complexity:
\[
m_{\mathrm{CoT}}\text{ can be independent of }T,
\qquad
R_{\mathrm{reason}}=\Omega(n)\text{ on BAPO-hard tasks.}
\]

\section{겉으로 모순처럼 보이는 문장 정리}

\begin{enumerate}
\item ``긴 CoT는 계산 능력을 늘린다.'' --- \textbf{가능}. Expressivity 관점이다.
\item ``긴 CoT는 더 많은 training sample을 요구한다.'' --- \textbf{필수 아님}. CoT supervision에서는 $T$-independent bound가 가능하다.
\item ``CoT는 sample complexity를 줄인다.'' --- \textbf{조건부}. Intermediate trace가 informative해야 한다.
\item ``reasoning token은 압축하면 된다.'' --- \textbf{항상 아님}. 일부 task에는 token lower bound가 있다.
\item ``Transformer가 GD/backprop을 구현할 수 있으니 실제 SGD가 그 circuit을 찾는다.'' --- \textbf{따라오지 않음}. Existence와 optimization은 별도다.
\item ``CoT가 유용하니 설명도 faithful하다.'' --- \textbf{따라오지 않음}. Faithfulness는 별도 성질이다.
\end{enumerate}

\section{발표에서 예상되는 질문}

\subsection*{Q1. Week 5의 computational theory와 이번 주가 겹치지 않나?}
Week 5의 목적은 Transformer 계산 모델의 upper/lower class를 이해하는 것이었다.
이번 주는 CoT를 중심으로 \textbf{계산 자원 + forward-pass 내부 optimization + 학습 복잡도 + token lower bound}를 연결한다. 특히 ICL--GD/backprop construction과 Joshi/Hanneke/Altabaa/BAPO가 새 연결점이다.

\subsection*{Q2. CoT sample complexity가 $T$에 독립이면 긴 trace가 공짜인가?}
아니다. Training example 수가 $T$에 독립일 수 있다는 뜻이지, 한 sample 안의 token 수나 training FLOPs가 독립이라는 뜻이 아니다.

\subsection*{Q3. BAPO lower bound가 실제 GPT류 LLM에 그대로 적용되나?}
아니다. BAPO는 bounded information flow를 추상화한 모델이다.
실제 model experiment와 정성적 일치는 보고되지만 theorem의 직접 적용이라고 말하면 과장이다.

\subsection*{Q4. ICL--GD 동치성이 실제 LLM도 backprop을 한다는 뜻인가?}
아니다. von Oswald/Bai/Wu의 핵심 정리는 특정 architecture에 대한 explicit construction이다. 일부 단순 regression 실험에서는 학습된 weight와 construction의 정렬 evidence가 있지만, 자연어 LLM 전체에 대한 mechanistic identity는 아니다.

\subsection*{Q5. CoT information이 크면 무조건 좋은 reasoning trace인가?}
통계적 hypothesis discrimination 측면에서는 유리하다.
하지만 human interpretability, factuality, causal faithfulness를 직접 측정하는 양은 아니다.

\begin{presentbox}
이번 주 전체를 한 문장으로 정리하면 이렇습니다. Transformer는 한 forward pass 안에 optimizer/backprop circuit을 구현할 수 있고, autoregressive CoT는 그 forward pass를 반복해 serial compute를 추가합니다. Training trace는 hypothesis를 더 빨리 식별하게 만들 수 있지만, inference reasoning token에는 별도의 task-dependent lower bound가 있습니다.
\end{presentbox}

% =====================================================================
\chapter{발표용 압축 정리}
% =====================================================================

\section{10분 설명 순서}

\begin{enumerate}
\item \textbf{1분}: CoT를 $z_{t+1}=f(x,z_{1:t})$라는 repeated computation으로 정의.
\item \textbf{2분}: Merrill--Sabharwal --- decoding step이 serial computation resource라는 점.
\item \textbf{3분}: ICL $\leftrightarrow$ optimization --- linear attention의 1-step GD, Bai의 algorithm selection, Wu의 $(2N+4)L$ backprop simulation.
\item \textbf{2분}: Joshi/Hanneke/Altabaa --- intermediate trace가 sample complexity와 hypothesis discrimination을 어떻게 바꾸는가.
\item \textbf{2분}: BAPO + caveat --- inference token lower bound, existence vs learned mechanism, faithfulness 구분.
\end{enumerate}

\section{칠판에 남길 식}

\[
z_{t+1}=f(xz_{1:t}),
\qquad
\TIME(t(n))\subseteq\CoT(t(n)).
\]
\[
\Delta W
=-\frac{\eta}{N}\sum_i(Wx_i-y_i)x_i^\top
\quad\Longleftrightarrow\quad
\text{one linear-attention GD update}.
\]
\[
\begin{gathered}
L\text{ GD steps on an }N\text{-layer network}\\
\Longleftarrow (2N+4)L\text{ Transformer layers}.
\end{gathered}
\]
\[
\begin{gathered}
m^{\mathrm{CoT},T}\text{ can be independent of }T,\qquad
m\sim \frac{d}{\Icot(\epsilon)},\\
R_{\mathrm{reason}}=\Omega(n)\quad\text{on BAPO-hard tasks}.
\end{gathered}
\]

\section{마지막에 반드시 붙일 caveat}

\begin{warningbox}
이번 주의 theorem들은 서로 다른 formal model을 사용한다. Merrill--Sabharwal의 decoder construction, von Oswald/Bai/Wu의 in-context optimizer constructions, Joshi/Hanneke의 PAC model, Altabaa의 CoT hypothesis class, Tomlinson et al.의 BAPO abstraction을 하나의 universal theorem처럼 합치면 안 된다. 특히 \textbf{optimizer를 구현하는 weight가 존재한다}와 \textbf{SGD가 실제로 그 weight를 찾는다}를 구분해야 한다.
\end{warningbox}

\begin{presentbox}
CoT 이론에서 중요한 구분은 네 가지입니다. 몇 decoding step을 쓰면 무엇을 계산할 수 있는가, 한 forward pass 안에 어떤 optimizer를 구현할 수 있는가, 그 rule을 몇 sample로 학습할 수 있는가, 그리고 실제 문제를 풀 때 최소 몇 reasoning token이 필요한가입니다. 이 네 질문을 분리하면 ICL과 CoT 결과들이 한 그림 안에서 충돌 없이 정리됩니다.
\end{presentbox}

% =====================================================================
% Bibliography (thebibliography prints the single heading and TOC entry)
% =====================================================================

\begin{thebibliography}{99}

\bibitem{merrill2024}
William Merrill and Ashish Sabharwal.
\newblock The Expressive Power of Transformers with Chain of Thought.
\newblock International Conference on Learning Representations (ICLR), 2024.

\bibitem{joshi2025}
Nirmit Joshi, Gal Vardi, Adam Block, Surbhi Goel, Zhiyuan Li, Theodor Misiakiewicz, and Nathan Srebro.
\newblock A Theory of Learning with Autoregressive Chain of Thought.
\newblock Proceedings of the 38th Conference on Learning Theory (COLT), PMLR 291:3161--3212, 2025.

\bibitem{altabaa2025}
Awni Altabaa, Omar Montasser, and John D. Lafferty.
\newblock CoT Information: Improved Sample Complexity under Chain-of-Thought Supervision.
\newblock Advances in Neural Information Processing Systems (NeurIPS), 2025.

\bibitem{hanneke2026}
Steve Hanneke, Idan Mehalel, and Shay Moran.
\newblock Sample Complexity of Autoregressive Reasoning: Chain-of-Thought vs. End-to-End.
\newblock arXiv:2604.12013, 2026.

\bibitem{tomlinson2026}
Kiran Tomlinson, Tobias Schnabel, Adith Swaminathan, and Jennifer Neville.
\newblock Reasoning about Reasoning: BAPO Bounds on Chain-of-Thought Token Complexity in LLMs.
\newblock International Conference on Machine Learning (ICML), 2026.

\bibitem{vonoswald2023}
Johannes von Oswald, Eyvind Niklasson, Ettore Randazzo, Jo\~ao Sacramento, Alexander Mordvintsev, Andrey Zhmoginov, and Max Vladymyrov.
\newblock Transformers Learn In-Context by Gradient Descent.
\newblock Proceedings of the 40th International Conference on Machine Learning (ICML), PMLR 202, 2023. arXiv:2212.07677.

\bibitem{akyurek2023}
Ekin Aky\"urek, Dale Schuurmans, Jacob Andreas, Tengyu Ma, and Denny Zhou.
\newblock What Learning Algorithm Is In-Context Learning? Investigations with Linear Models.
\newblock International Conference on Learning Representations (ICLR), 2023. arXiv:2211.15661.

\bibitem{bai2023}
Yu Bai, Fan Chen, Huan Wang, Caiming Xiong, and Song Mei.
\newblock Transformers as Statisticians: Provable In-Context Learning with In-Context Algorithm Selection.
\newblock arXiv:2306.04637, 2023.

\bibitem{wu2025}
Weimin Wu, Maojiang Su, Jerry Yao-Chieh Hu, Zhao Song, and Han Liu.
\newblock In-Context Deep Learning via Transformer Models.
\newblock Proceedings of the 42nd International Conference on Machine Learning (ICML), PMLR 267, 2025. arXiv:2411.16549.

\bibitem{merrill2023}
William Merrill and Ashish Sabharwal.
\newblock The Parallelism Tradeoff: Limitations of Log-Precision Transformers.
\newblock Transactions of the Association for Computational Linguistics, 2023.

\end{thebibliography}

\end{document}
